\subsection{Construction statements of work}

\paragraph{concrete short.}
\begin{quote}\small Summary Statement of Work, Trade: Concrete, Section 03 30 00 Cast-in-Place Concrete. Scope: 1. Furnish and place 185 CY of 4,000 psi normal-weight concrete for spread footings and grade beams, 4 in slump, 6 percent air. 2. Furnish and install 9.2 tons of ASTM A615 Grade 60 reinforcing steel with 3 in cover at earth-formed surfaces. 3. Provide edge formwork, 7-day moist curing, and one set of four ASTM C31 test cylinders per 50 CY placed, broken per ASTM C39 at 7 and 28 days. Major exclusions: excavation, backfill, and dewatering by the Sitework contractor.\end{quote}

\paragraph{concrete long.}
\begin{quote}\small Statement of Work, Trade: Concrete (Cast-in-Place), Building B Warehouse Expansion. Spec sections used: 03 20 00 Concrete Reinforcing, 03 30 00 Cast-in-Place Concrete, 03 35 00 Concrete Finishing, 07 26 00 Vapor Retarders. The Concrete subcontractor shall furnish all labor, materials, equipment, and supervision for the following scope. Foundations: 1. Spread footings, 24 each, 6 ft x 6 ft x 18 in, 4,000 psi concrete, 48 CY total. 2. Continuous wall footings, 620 LF, 2 ft wide x 12 in deep, 4,000 psi, 46 CY. 3. Foundation walls, 620 LF, 8 in thick x 4 ft high, 4,000 psi, 61 CY, formed both faces. 4. Loading dock pits, 4 each, 10 ft x 8 ft x 4 ft deep with 8 in walls and slab, 5,000 psi, 22 CY, with embedded steel angles at all edges. 5. Anchor bolts: install 96 each 1 in diameter x 24 in ASTM F1554 Grade 36 anchor bolts per templates furnished by the Steel Erector. Slabs: 6. Slab-on-grade, 42,000 SF, 6 in thick, 4,500 psi, 778 CY, placed over 15-mil vapor retarder per ASTM E1745 Class A. 7. Equipment pads, 6 each, 4 ft x 6 ft x 6 in, 4,000 psi, 3 CY, with 3/4 in chamfered edges. 8. Sawcut control joints at 15 ft maximum spacing each way, 1/4 of slab depth, cut within 12 hours of finishing. 9. Slab finish: hard steel trowel, floor flatness FF 35 and floor levelness FL 25 per ASTM E1155, tested within 72 hours of placement. Reinforcing: 10. Slab reinforcing: \#4 bars at 18 in on center each way, ASTM A615 Grade 60, on chairs at 3 ft on center. 11. Footing and wall reinforcing: 14.5 tons ASTM A615 Grade 60, 3 in cover against earth, 2 in cover at formed surfaces; furnish mill certificates. Materials and quality: 12. Mix designs: maximum water-cement ratio 0.45, 4 in slump plus or minus 1 in, air entrainment 6 percent plus or minus 1.5 percent for exterior concrete, Type II cement, 25 percent Class F fly ash replacement permitted. 13. Curing: dissipating-resin curing compound at 200 SF per gallon within 30 minutes of final finishing; moist-cure footings and walls for 7 days. 14. Cold weather: maintain concrete at 55 degrees F minimum for 3 days after placement when ambient temperature is below 40 degrees F, per ACI 306. 15. Tolerances per ACI 117: footings plus 2 in or minus 1/2 in in plan; walls plus or minus 1/4 in in 10 ft. 16. Testing: one set of four cylinders per 100 CY or fraction thereof per day per mix, ASTM C31 field cured, ASTM C39 breaks at 7 and 28 days, slump and air per ASTM C143 and ASTM C231 at each set. 17. Submittals within 14 calendar days of award: concrete mix designs, reinforcing shop drawings, vapor retarder and curing compound product data, mill certificates. Schedule: 18. Footings complete within 15 working days of mobilization; slab-on-grade placed in 6 pours of approximately 7,000 SF each; 28-day design strength required before steel erection loads. Exclusions (by others): 19. Excavation, backfill, dewatering, and compacted subgrade by the Sitework contractor. 20. The 4 in granular capillary break under slabs by the Sitework contractor. 21. Testing agency fees paid by the Owner; anchor bolt templates furnished by the Steel Erector; dumpster for concrete debris provided by the General Contractor. Alternates: Alternate 1: substitute 6 x 6 W2.9 x W2.9 welded wire reinforcement for the \#4 slab bars, deduct \$18,400. Base bid: \$1,284,600 including a 5 percent contingency.\end{quote}

\subsection{Sixteen-domain breadth prompts}
Each prompt is followed by its ten anchor facts (type: text).

\paragraph{recipe.}
\begin{quote}\small Chicken tikka masala for 4 people. Cut 1.5 lb boneless chicken thighs into 1-inch pieces and marinate 30 minutes in 1 cup plain yogurt, 2 tablespoons lemon juice, 2 teaspoons garam masala, 1 teaspoon turmeric, and 1 teaspoon salt. Heat 2 tablespoons ghee in a large skillet over medium-high heat and sear the chicken 4 minutes per side until charred at the edges, then set aside. In the same pan cook 1 diced onion for 5 minutes, add 4 minced garlic cloves and 1 tablespoon grated ginger, and cook 1 minute more. Stir in 2 tablespoons tomato paste and 1 teaspoon smoked paprika, then pour in one 14-oz can of crushed tomatoes and simmer 10 minutes. Add 3/4 cup heavy cream and the chicken, and simmer 8 minutes until the sauce thickens. Finish with 2 tablespoons chopped cilantro. Serve over basmati rice with warm naan.\end{quote}
\noindent\footnotesize Anchors: num: ``1.5 lb''; num: ``30 minutes''; num: ``2 teaspoons garam masala''; num: ``1 teaspoon turmeric''; num: ``4 minutes per side''; num: ``4 minced garlic cloves''; num: ``14-oz''; num: ``10 minutes''; num: ``3/4 cup heavy cream''; num: ``8 minutes''\normalsize

\paragraph{lab protocol.}
\begin{quote}\small Protocol: plasmid miniprep by alkaline lysis. Pellet 3 mL of overnight E. coli culture at 8,000 x g for 2 minutes and discard the supernatant. Resuspend the pellet in 250 uL of Buffer P1 containing RNase A at 100 ug/mL and vortex until no clumps remain. Add 250 uL of Buffer P2 (200 mM NaOH, 1 percent SDS), invert the tube 6 times, and incubate no longer than 5 minutes at room temperature. Neutralise with 350 uL of Buffer N3, invert 6 times, and centrifuge at 13,000 rpm for 10 minutes. Transfer the cleared lysate to a silica spin column, spin 60 seconds, and discard the flow-through. Wash with 500 uL of Buffer PB and then 750 uL of Buffer PE, spinning 60 seconds after each wash and once more dry for 1 minute. Elute in 50 uL of Buffer EB warmed to 50 degrees C, letting it stand 1 minute before the final spin. Expected yield is 5 to 15 ug; verify by measuring A260/A280, which should be between 1.8 and 2.0.\end{quote}
\noindent\footnotesize Anchors: num: ``3 mL''; num: ``8,000 x g''; num: ``250 uL''; num: ``100 ug/mL''; num: ``200 mM NaOH''; num: ``350 uL''; num: ``13,000 rpm''; num: ``750 uL''; num: ``50 degrees C''; num: ``1.8 and 2.0''\normalsize

\paragraph{lockout tagout.}
\begin{quote}\small Lockout/tagout procedure LOTO-114 for Hydraulic Press 7 in Building C. Only authorised employees trained under OSHA 29 CFR 1910.147 may perform this procedure. Step 1: Notify the shift supervisor at extension 4471 and all affected employees that the press will be shut down. Step 2: Stop the press using the operator stop button and confirm the ram is fully retracted. Step 3: Open the main electrical disconnect MCC-3, breaker 12, and apply a red personal lock and tag. Step 4: Close the hydraulic supply valve HV-7A and apply a second lock. Step 5: Bleed residual pressure at the gauge port until the gauge reads 0 psi; the accumulator holds up to 3,000 psi and must be fully vented. Step 6: Attempt a restart to verify zero energy state, then return the controls to off. Step 7: Perform the work. Step 8: Remove tools, replace guards, confirm all personnel are clear, and remove locks in reverse order. Each lock stays with its owner; a lock may be cut only by the plant manager with a written release. Review this procedure annually, next due 15 March 2027.\end{quote}
\noindent\footnotesize Anchors: id: ``LOTO-114''; name: ``Hydraulic Press 7''; id: ``29 CFR 1910.147''; num: ``extension 4471''; id: ``MCC-3''; num: ``breaker 12''; id: ``HV-7A''; num: ``3,000 psi''; name: ``plant manager''; date: ``15 March 2027''\normalsize

\paragraph{concrete sow.}
\begin{quote}\small Statement of Work, Trade: Concrete, Section 03 30 00 Cast-in-Place Concrete, Building B Warehouse Expansion. Scope: 1. Spread footings, 24 each, 6 ft x 6 ft x 18 in, 4,000 psi concrete, 48 CY total. 2. Foundation walls, 620 LF, 8 in thick x 4 ft high, 4,000 psi, 61 CY, formed both faces. 3. Slab-on-grade, 42,000 SF, 6 in thick, 4,500 psi, 778 CY, over a 15-mil vapor retarder per ASTM E1745. 4. Reinforcing: 14.5 tons of ASTM A615 Grade 60 steel with 3 in cover against earth; slab bars \#4 at 18 in on center each way. 5. Mix design: maximum water-cement ratio 0.45, 4 in slump, 6 percent air for exterior concrete. 6. Sawcut control joints at 15 ft maximum spacing within 12 hours of finishing. 7. Testing: one set of four ASTM C31 cylinders per 100 CY, broken at 7 and 28 days. 8. Submittals within 14 calendar days of award. Exclusions: excavation, backfill, and dewatering by the Sitework contractor; anchor bolt templates by the Steel Erector. Base bid: \$1,284,600 including a 5 percent contingency.\end{quote}
\noindent\footnotesize Anchors: id: ``03 30 00''; num: ``24 each''; num: ``4,000 psi''; num: ``620 LF''; num: ``42,000 SF''; num: ``778 CY''; id: ``ASTM E1745''; num: ``14.5 tons''; num: ``0.45''; num: ``\$1,284,600''\normalsize

\paragraph{construction rfi.}
\begin{quote}\small Request for Information RFI-047, Project: Maple Street Elementary Addition, Date: 12 August 2026, From: Harlan Mechanical, To: Kessler Architects, Response required by 19 August 2026. Subject: Rooftop unit RTU-3 curb height conflict. Drawing M-401 shows RTU-3 on a 14-inch roof curb at grid line D-7, but structural drawing S-201 detail 5 shows a 24-inch curb to clear the tapered insulation, which is 10 inches thick at that location per A-501. Specification Section 23 74 13 paragraph 2.4.B requires the curb to be factory-fabricated and shipped with the unit, and the unit was released to the manufacturer on 30 July 2026 with the 14-inch curb. Question 1: Which curb height governs? Question 2: If 24 inches is required, will the Owner accept a field-fabricated 12-inch extension welded to the factory curb, or must the unit be re-released at a cost of \$6,850 and 6 weeks of delay? Question 3: Does the condensate drain routing on M-401 change if the curb height changes? Contractor recommendation: accept the field extension with flashing per SMACNA Figure 8-2. Please respond in writing; work at grid D-7 is on hold.\end{quote}
\noindent\footnotesize Anchors: id: ``RFI-047''; name: ``Harlan Mechanical''; name: ``Kessler Architects''; date: ``19 August 2026''; id: ``RTU-3''; id: ``M-401''; id: ``D-7''; num: ``14-inch''; num: ``24-inch''; num: ``\$6,850''\normalsize

\paragraph{change order.}
\begin{quote}\small Change Order Request COR-019, Project: Riverside Clinic Renovation, Contract No. 2026-114, Date: 3 September 2026. Reason: unforeseen asbestos-containing floor mastic discovered under the existing vinyl tile in Rooms 112 through 118, confirmed by laboratory report EnvTest-88213 dated 28 August 2026. Scope: abatement of 2,340 square feet of mastic by a licensed abatement subcontractor, air monitoring, and disposal at a permitted landfill. Cost breakdown: abatement labor, 4 workers for 6 days at \$68 per hour, \$13,056; abatement materials and encapsulant, \$2,410; air monitoring and clearance testing, \$3,200; disposal, 3 lined containers at \$850 each, \$2,550; subcontractor overhead and profit at 15 percent, \$3,182; general contractor markup at 10 percent, \$2,440. Total this change order: \$26,838. Schedule impact: 8 working days added to the Substantial Completion date, moving it from 14 November 2026 to 26 November 2026. This proposal is valid for 30 days. Owner approval is required before abatement begins; the flooring installer has been placed on standby.\end{quote}
\noindent\footnotesize Anchors: id: ``COR-019''; id: ``2026-114''; num: ``2,340 square feet''; id: ``EnvTest-88213''; num: ``\$68 per hour''; num: ``\$13,056''; num: ``\$3,200''; num: ``15 percent''; num: ``\$26,838''; date: ``26 November 2026''\normalsize

\paragraph{contract clause.}
\begin{quote}\small Article 11, Indemnification. To the fullest extent permitted by law, the Subcontractor shall indemnify, defend, and hold harmless the Contractor, the Owner, and their respective officers, agents, and employees from and against all claims, damages, losses, and expenses, including reasonable attorneys' fees, arising out of or resulting from the performance of the Subcontract Work, but only to the extent caused by the negligent acts or omissions of the Subcontractor, its sub-subcontractors, or anyone for whose acts they may be liable. This obligation shall not be construed to negate or reduce any other right of indemnity that would otherwise exist. The Subcontractor shall maintain commercial general liability insurance of not less than \$2,000,000 per occurrence and \$4,000,000 aggregate, naming the Contractor and Owner as additional insureds on a primary and non-contributory basis, with a waiver of subrogation. Certificates of insurance shall be delivered within 10 days of execution and not less than 30 days before any policy expires. Notice of any claim shall be given in writing within 14 days of the Subcontractor first becoming aware of it. This Article survives termination of the Subcontract for a period of 3 years.\end{quote}
\noindent\footnotesize Anchors: name: ``Article 11''; name: ``Subcontractor''; name: ``Owner''; num: ``\$2,000,000''; num: ``\$4,000,000''; name: ``additional insureds''; name: ``waiver of subrogation''; num: ``10 days''; num: ``14 days''; num: ``3 years''\normalsize

\paragraph{building code.}
\begin{quote}\small Section 1005, Means of Egress Sizing (excerpt). 1005.3.1 Stairways: the capacity of stairways shall be calculated at 0.3 inch of width per occupant, except that in buildings equipped throughout with an automatic sprinkler system and an emergency voice alarm system, 0.2 inch per occupant is permitted. 1005.3.2 Other egress components: 0.2 inch per occupant, reduced to 0.15 inch where the sprinkler and voice alarm exception applies. 1005.7 Encroachment: doors opening into the path of egress travel shall not reduce the required width by more than 7 inches when fully open, and by not more than one-half at any point during the swing. Minimum stair width is 44 inches where the occupant load served exceeds 50, and 36 inches otherwise. Handrails may project not more than 4.5 inches into the required width on each side. Corridors serving an occupant load of 50 or more shall be not less than 44 inches wide; in Group E occupancies with more than 100 occupants, 72 inches. Where two exits are required, each shall be sized for not less than 50 percent of the total occupant load.\end{quote}
\noindent\footnotesize Anchors: id: ``1005.3.1''; num: ``0.3 inch''; num: ``0.2 inch''; num: ``0.15 inch''; num: ``7 inches''; num: ``44 inches''; num: ``36 inches''; num: ``4.5 inches''; name: ``Group E''; num: ``72 inches''\normalsize

\paragraph{discharge instructions.}
\begin{quote}\small Discharge instructions for Robert Alvarez, age 67, discharged 4 September 2026 after treatment for community-acquired pneumonia. Medications: amoxicillin-clavulanate 875 mg by mouth every 12 hours for 7 more days, take with food; azithromycin 250 mg once daily for 4 more days; continue lisinopril 10 mg every morning; acetaminophen 650 mg every 6 hours as needed for fever, not to exceed 3,000 mg in 24 hours. Do not take ibuprofen while on lisinopril without asking your doctor. Use the incentive spirometer 10 breaths every hour while awake. Drink at least 2 liters of fluid daily unless told otherwise. Call the clinic at 555-0142 or go to the emergency department if your temperature exceeds 101.5 degrees F, if you become short of breath at rest, or if you cough up blood. Follow-up appointment with Dr. Priya Natarajan on 11 September 2026 at 9:30 AM; bring this sheet and all medication bottles. A repeat chest X-ray is scheduled for 6 weeks after discharge to confirm the infiltrate has cleared. No driving for 48 hours after the last dose of the cough medicine containing codeine.\end{quote}
\noindent\footnotesize Anchors: name: ``Robert Alvarez''; num: ``875 mg''; num: ``every 12 hours''; num: ``250 mg''; num: ``10 mg''; num: ``3,000 mg''; num: ``555-0142''; num: ``101.5 degrees F''; name: ``Dr. Priya Natarajan''; date: ``11 September 2026''\normalsize

\paragraph{python code.}
\begin{quote}\small Python module for rate limiting. def token\_bucket(capacity: int = 100, refill\_rate: float = 10.0): creates a closure holding two variables, tokens initialised to capacity and last\_refill initialised to time.monotonic(). The inner function allow(cost: int = 1) -> bool first computes elapsed = time.monotonic() - last\_refill, then sets tokens = min(capacity, tokens + elapsed * refill\_rate) and updates last\_refill. If tokens >= cost it subtracts cost and returns True; otherwise it returns False without changing tokens. The function raises ValueError if capacity <= 0 or refill\_rate <= 0. A second helper, retry\_after(cost: int) -> float, returns max(0.0, (cost - tokens) / refill\_rate), the seconds until enough tokens will exist. Both helpers are returned as a tuple (allow, retry\_after). Usage: allow, retry\_after = token\_bucket(capacity=50, refill\_rate=5.0); if not allow(3): sleep(retry\_after(3)). The module depends only on the standard library time module and is thread-unsafe by design; callers must wrap it in a threading.Lock when sharing across threads. Version 1.4.2, licensed under the MIT License.\end{quote}
\noindent\footnotesize Anchors: id: ``token\_bucket''; num: ``capacity: int = 100''; num: ``refill\_rate: float = 10.0''; id: ``time.monotonic()''; id: ``allow(cost: int = 1)''; id: ``ValueError''; id: ``retry\_after''; id: ``threading.Lock''; num: ``1.4.2''; name: ``MIT License''\normalsize

\paragraph{api reference.}
\begin{quote}\small Endpoint: POST /v2/invoices. Creates an invoice for a customer. Authentication: Bearer token with the invoices:write scope. Request body (JSON): customer\_id (string, required, format cus\_ followed by 14 characters); currency (string, required, ISO 4217 code such as USD); line\_items (array, required, 1 to 250 objects each with description, quantity as an integer, and unit\_amount in minor units); due\_date (string, optional, ISO 8601 date, defaults to 30 days after creation); metadata (object, optional, up to 50 keys). Response 201 Created returns the invoice object with id (format inv\_ plus 14 characters), status set to draft, and total computed as the sum of quantity times unit\_amount. Errors: 400 invalid\_request when a required field is missing; 402 payment\_required when the account is past due; 404 not\_found when customer\_id does not exist; 409 conflict when an Idempotency-Key header was already used with a different body; 429 rate\_limited when more than 100 requests per minute are sent. Idempotency keys are retained for 24 hours. This endpoint was introduced in API version 2025-11-01 and replaces POST /v1/invoices, which is deprecated and will be removed on 1 June 2027.\end{quote}
\noindent\footnotesize Anchors: id: ``POST /v2/invoices''; id: ``invoices:write''; id: ``customer\_id''; id: ``ISO 4217''; num: ``250''; num: ``201 Created''; num: ``402''; num: ``409''; num: ``24 hours''; date: ``1 June 2027''\normalsize

\paragraph{scientific abstract.}
\begin{quote}\small Background: Whether a 12-week structured walking programme reduces blood pressure in adults with stage 1 hypertension remains uncertain. Methods: We randomised 248 participants aged 45 to 70 years at 6 primary care clinics to either supervised walking for 40 minutes 5 days per week or usual care. The primary outcome was change in systolic blood pressure at 12 weeks. Results: 231 participants (93 percent) completed follow-up. Mean systolic pressure fell by 7.8 mmHg in the walking group versus 2.1 mmHg with usual care, a between-group difference of 5.7 mmHg (95 percent confidence interval 3.4 to 8.0; p < 0.001). Diastolic pressure fell by 3.2 mmHg more in the walking group (p = 0.004). Adherence averaged 4.3 sessions per week. Two participants in the walking group reported musculoskeletal injuries; no serious adverse events occurred. The effect was larger in participants with a baseline body mass index above 30 (interaction p = 0.03). Conclusions: A supervised walking programme produced a clinically meaningful reduction in systolic blood pressure over 12 weeks. Trial registration NCT05873310.\end{quote}
\noindent\footnotesize Anchors: num: ``12-week''; num: ``248''; num: ``45 to 70''; num: ``40 minutes''; num: ``93 percent''; num: ``7.8 mmHg''; num: ``5.7 mmHg''; num: ``p < 0.001''; num: ``4.3 sessions''; id: ``NCT05873310''\normalsize

\paragraph{financial summary.}
\begin{quote}\small Northwind Logistics Inc. reported results for the second quarter ended 30 June 2026. Revenue was \$412.6 million, up 9.4 percent from \$377.1 million in the prior-year quarter, driven by a 14 percent increase in cross-border freight volume. Gross margin expanded to 23.8 percent from 21.9 percent as fuel costs declined. Operating income was \$38.2 million, and net income was \$27.5 million, or \$1.12 per diluted share, compared with \$0.89 a year earlier. Adjusted EBITDA reached \$61.7 million. Operating cash flow for the six months was \$88.3 million; capital expenditures were \$31.0 million, mostly for 140 new tractors. The company ended the quarter with \$96.4 million in cash and \$210 million of undrawn revolver capacity, and repurchased 1.1 million shares for \$52 million. Management raised full-year revenue guidance to a range of \$1.62 billion to \$1.66 billion and reaffirmed an expected effective tax rate of 24 percent. The board declared a quarterly dividend of \$0.18 per share payable on 15 October 2026 to holders of record on 30 September 2026.\end{quote}
\noindent\footnotesize Anchors: name: ``Northwind Logistics''; num: ``\$412.6 million''; num: ``9.4 percent''; num: ``23.8 percent''; num: ``\$27.5 million''; num: ``\$1.12''; num: ``\$61.7 million''; num: ``1.1 million shares''; num: ``\$0.18''; date: ``15 October 2026''\normalsize

\paragraph{news report.}
\begin{quote}\small PORTLAND, Oregon, 5 September 2026. A 6.1-magnitude earthquake struck off the Oregon coast at 4:47 a.m. local time on Friday, the U.S. Geological Survey said, centred about 62 miles west of Newport at a depth of 10 kilometres. No tsunami warning was issued. Oregon Emergency Management reported minor damage to a bridge on U.S. Route 20 near Toledo and a power outage affecting roughly 8,400 customers in Lincoln County, most restored by 9 a.m. Two people were treated for minor injuries at Samaritan Pacific Communities Hospital, a spokesperson said. Governor Elena Whitcombe said state inspectors would examine 14 coastal bridges over the weekend. Amtrak suspended Cascades service between Portland and Eugene for 3 hours while tracks were inspected. Seismologist Daniel Okafor of Oregon State University said the quake occurred on the Cascadia subduction zone's outer fault system and that aftershocks above magnitude 4 were possible for several days. The last quake of comparable size in the region was a 6.3 event in 2004.\end{quote}
\noindent\footnotesize Anchors: num: ``6.1''; num: ``4:47 a.m.''; name: ``Newport''; num: ``10 kilometres''; id: ``U.S. Route 20''; num: ``8,400''; name: ``Elena Whitcombe''; num: ``14 coastal bridges''; name: ``Daniel Okafor''; num: ``2004''\normalsize

\paragraph{travel itinerary.}
\begin{quote}\small Itinerary for Maria Chen, confirmation code QX7R4M. Monday 21 September 2026: depart Boston Logan Terminal B on JetBlue flight B6 1523 at 7:15 AM, arrive Chicago O'Hare Terminal 5 at 9:02 AM, seat 14C. Pick up a mid-size rental car from Hertz, reservation H8823391, at the O'Hare rental center. Check in after 3:00 PM at the Kimpton Gray Hotel, 122 West Monroe Street, Chicago, reservation 44192077, 2 nights, king room, \$289 per night. Tuesday 22 September: client meeting at 10:00 AM with Ardent Partners, 233 South Wacker Drive, Suite 4100, contact Luis Ortega, 312-555-0177; dinner reservation at 7:30 PM for 4 at Girl \& the Goat. Wednesday 23 September: return the car by 12:00 PM, depart O'Hare on B6 1530 at 2:40 PM, arrive Boston at 6:05 PM. Total airfare \$486.40 charged to the corporate card ending 7731. Travel insurance policy TI-90315 covers trip cancellation up to \$2,500.\end{quote}
\noindent\footnotesize Anchors: name: ``Maria Chen''; id: ``QX7R4M''; id: ``B6 1523''; num: ``7:15 AM''; num: ``14C''; id: ``H8823391''; name: ``122 West Monroe Street''; id: ``44192077''; name: ``Luis Ortega''; num: ``\$486.40''\normalsize

\paragraph{literary prose.}
\begin{quote}\small The lighthouse keeper's daughter kept a ledger of the fog. Not the weather service kind, with its dutiful columns of visibility and wind, but a private accounting of how the grey came in: some mornings like a tide, filling the harbour from the bottom up until the masts of the fishing boats stood in it like reeds; other mornings all at once, a door closing on the world. Her father said fog was only water that had forgotten which way was down. She wrote that in the ledger too. In October the fog stayed nine days without lifting, and she learned the sound of the town without its shapes: the bakery shutter at six, the school bell that seemed to ring from underwater, the horn on the point answering itself every thirty seconds like a man who has lost the thread of an argument. When it finally cleared she found the harbour exactly as it had been, and was surprised to feel disappointed, as though the fog had promised to take something with it when it went and had failed to keep its word.\end{quote}
\noindent\footnotesize Anchors: name: ``lighthouse keeper's daughter''; name: ``ledger''; name: ``harbour''; name: ``fishing boats''; name: ``forgotten which way was down''; name: ``October''; num: ``nine days''; name: ``bakery shutter''; name: ``school bell''; num: ``thirty seconds''\normalsize
